% ============================================================
%  Lecture 2 — Knapsack, complexity, RCPSP, shop scheduling
%  (still Chapter 1: Introduction to Scheduling)
% ============================================================

\section{The Knapsack Problem}
\label{sec:knapsack}

Before continuing with scheduling problems, we take a detour
through a simpler combinatorial problem that lets us sharpen
three important ideas: greedy heuristics, upper bounds, and the
distinction between easy and hard problems.%
\aside{Not a scheduling problem, but the
perfect sandbox for greedy, bounds, and
complexity concepts.}

\begin{definition}[0/1 Knapsack Problem]
\label{def:knapsack}
We are given $n$ items, each with a \emph{weight} $w_i$ and a
\emph{value} $v_i$, and a knapsack of capacity~$W$.
We must choose a subset of items that fits in the knapsack
and maximises total value:
\[
  \max \sum_{i=1}^{n} v_i\, x_i
  \qquad\text{s.t.}\quad
  \sum_{i=1}^{n} w_i\, x_i \le W,
  \qquad x_i \in \{0,1\}.
\]
\end{definition}

In plain words: every item is either taken ($x_i = 1$) or
left behind ($x_i = 0$).  We want as much value as possible
without exceeding the weight limit.  Think of packing a
backpack for a hike, or---in computing---selecting files to
store on a disk with limited capacity.

% ------------------------------------------------------------
\subsection{Three greedy strategies}
\label{ssec:knapsack-greedy}
% ------------------------------------------------------------

\begin{example}[Six-item knapsack instance]
\label{ex:knapsack-six}
Consider six items with capacity $W = 100$:

\begin{figure}[H]
\centering
\begin{tabular}{lcccccc}
\toprule
\textbf{Item}   & 1   & 2   & 3   & 4  & 5  & 6 \\
\midrule
Weight $w_i$     & 100 & 50  & 45  & 20 & 10 & 5 \\
Value  $v_i$     & 40  & 35  & 18  & 4  & 10 & 2 \\
Ratio $v_i/w_i$  & 0.40& 0.70& 0.40& 0.20& 1.00& 0.40\\
\bottomrule
\end{tabular}
\caption{Data for the six-item knapsack instance
         (\cref{ex:knapsack-six}).}
\label{fig:knapsack-data}
\end{figure}
\end{example}

\noindent
Let's try three natural greedy rules on this instance:

\paragraph{Greedy 1: highest value first.}
Pick the item with the largest~$v_i$ that still fits.
\begin{itemize}[nosep]
  \item Select item~1 ($v = 40$, $w = 100$).
        Remaining capacity: $0$.
  \item No more items fit.
  \item \textbf{Total value: 40.}
\end{itemize}

\paragraph{Greedy 2: lightest first.}
Pick the item with the smallest~$w_i$ that still fits.
\begin{itemize}[nosep]
  \item Select item~6 ($w = 5$), then 5 ($w = 10$),
        then 4 ($w = 20$), then 3 ($w = 45$).
        Remaining capacity: $100 - 80 = 20$.
  \item Items~1 and~2 do not fit.
  \item \textbf{Total value: $2 + 10 + 4 + 18 = 34$.}
\end{itemize}

\paragraph{Greedy 3: best value-to-weight ratio.}
Pick the item with the highest $v_i / w_i$ that still fits.%
\aside{The classic knapsack heuristic.
Ties broken by largest item first.}
\begin{itemize}[nosep]
  \item Select item~5 (ratio 1.00, $w = 10$).
        Remaining: $90$.
  \item Select item~2 (ratio 0.70, $w = 50$).
        Remaining: $40$.
  \item Tied at ratio 0.40: items~1, 3, 6.
        Item~1 ($w = 100$) does not fit;
        item~3 ($w = 45$) does not fit;
        item~6 ($w = 5$) fits.  Remaining: $35$.
  \item Item~4 (ratio 0.20, $w = 20$) fits.  Remaining: $15$.
  \item \textbf{Total value: $10 + 35 + 2 + 4 = 51$.}
        Weight used: $85$.
\end{itemize}

\medskip\noindent
None of these finds the optimal solution.  By inspection, items
$\{2, 3, 6\}$ weigh exactly $50 + 45 + 5 = 100$ and achieve
value $35 + 18 + 2 = \mathbf{55}$.  The best greedy (ratio-based)
left 15 units of capacity unused, missing a better packing.

% ------------------------------------------------------------
\subsection{Upper bounds via relaxation}
\label{ssec:knapsack-bound}
% ------------------------------------------------------------

We found a solution worth~55 by inspection, but how do we
\emph{prove} no solution can do better?  The standard technique
is \emph{relaxation}: solve an easier version of the problem
whose optimal value is guaranteed to be at least as good.%
\aside{For a maximisation problem, a relaxation gives an
\emph{upper bound}.  For minimisation (like makespan), it gives
a \emph{lower bound}.}

\begin{definition}[Fractional Knapsack]
\label{def:fractional-knapsack}
The same formulation as \cref{def:knapsack}, but with
$x_i \in [0, 1]$ instead of $x_i \in \{0, 1\}$.
Items may be taken in fractions: if only half of item~$i$
fits, we take half and get $\tfrac{1}{2} v_i$.
\end{definition}

Since any 0/1 solution is also a valid fractional solution,
the fractional optimum is an upper bound on the integer optimum.
And the fractional problem has a simple exact greedy:

\begin{observation}
The greedy algorithm that sorts items by decreasing $v_i / w_i$
and takes as much of each item as fits is \emph{optimal}
for the fractional knapsack.
\end{observation}

\begin{example}[Fractional bound for \cref{ex:knapsack-six}]
\label{ex:fractional-bound}
Sort by ratio: item~5 (1.00), item~2 (0.70), then the
0.40-ratio group.
\begin{itemize}[nosep]
  \item Take all of item~5: weight 10, value 10. Remaining: 90.
  \item Take all of item~2: weight 50, value 35. Remaining: 40.
  \item Item~1 weighs 100; take $40/100 = 0.4$ of it:
        weight 40, value $0.4 \times 40 = 16$.
        Remaining: 0.
\end{itemize}
Total value: $10 + 35 + 16 = \mathbf{61}$.
Since the integer optimum $\le 61$ and we found a feasible
integer solution worth~55, the optimality gap is at most
$61 - 55 = 6$.
\end{example}

\begin{remark}
The fractional bound does not prove that 55 is optimal
(there could be an integer solution worth 56, 57, \ldots, 61).
For this small instance we verified optimality by inspection.
For large instances, tighter bounds or exact algorithms
(branch-and-bound, dynamic programming) are needed.
\end{remark}


% ============================================================
\section{Easy vs Hard Problems}
\label{sec:easy-hard}
% ============================================================

So far we have seen two \emph{easy} problems---project scheduling
and fractional knapsack---where a greedy algorithm finds the
optimum in polynomial time.  The 0/1 knapsack, by contrast,
is \emph{hard}: no known polynomial-time algorithm solves it
exactly.

\begin{definition}[$\NP$-hard (informal)]
\label{def:np-hard-informal}
A problem is $\NP$-hard if the best known exact algorithms
require time exponential in the input size.%
\aside{Whether $\NP$-hard problems truly
require exponential time is the open
$\mathsf{P} \stackrel{?}{=} \NP$ question.}
\end{definition}

The key structural difference: for easy problems, the algorithm
makes locally optimal choices with no branching.  For hard
problems, at each step there are multiple choices (which item
to take, which job to schedule next), and exploring all
combinations leads to an exponential number of candidate
solutions---roughly $2^n$ for $n$ binary decisions.

\begin{figure}[H]
\centering
\begin{tikzpicture}[
    level distance=14mm, sibling distance=28mm,
    edge from parent/.style={draw, -Stealth, thick},
    every node/.style={circle, draw=blue!60!black, fill=blue!8,
                       minimum size=6mm, font=\footnotesize},
    leaf/.style={fill=orange!15, draw=orange!60!black},
    chosen/.style={fill=violet!15, draw=violet!70!black},
    grayedge/.style={draw=gray!40, -Stealth},
    >=Stealth
  ]
  \node[chosen] {}
    child { node[chosen] {}
      child { node[chosen, leaf] {}
        edge from parent node[left, draw=none, fill=none,
          font=\scriptsize\itshape, text=violet!70!black] {A} }
      child { node[leaf] {}
        edge from parent[grayedge] }
      edge from parent node[left, draw=none, fill=none,
        font=\scriptsize\itshape, text=violet!70!black] {A} }
    child { node {}
      child { node[leaf] {}
        edge from parent[grayedge] }
      child { node[leaf] {}
        edge from parent[grayedge] }
      edge from parent[grayedge] };
\end{tikzpicture}
\caption{A greedy algorithm (violet path) picks one branch at
         each decision point, visiting $n$ nodes total.
         An exact algorithm may need to explore all $2^n$
         leaves.}
\label{fig:search-tree}
\end{figure}

\noindent
A \textbf{greedy} algorithm traverses a single root-to-leaf path
through this tree: fast ($n$~steps), but no guarantee of finding
the best leaf.
An \textbf{exact} algorithm explores enough of the tree to
certify optimality---potentially all of it.

\begin{observation}
Algorithms fall into two families:
\begin{itemize}[nosep]
  \item \textbf{Exact} algorithms guarantee the optimal
        solution, but may take exponential time on hard problems.
  \item \textbf{Heuristic} algorithms (greedy being the simplest
        subfamily) find good solutions quickly, without optimality
        guarantees.
\end{itemize}
This course covers both: we will implement greedy heuristics,
enumeration (backtracking), constraint programming (via MiniZinc),
and local search / metaheuristics.
\end{observation}


% ============================================================
\section{Resource-Constrained Project Scheduling (RCPSP)}
\label{sec:rcpsp}
% ============================================================

We now return to project scheduling and add the complication we
set aside earlier: \textbf{resources}.

\begin{definition}[Resource-Constrained Project Scheduling
  Problem---RCPSP]
\label{def:rcpsp}
We are given the same data as the basic project scheduling
problem (\cref{def:project-scheduling}), plus:
\begin{itemize}[nosep]
  \item $R$ types of \emph{renewable resources}, each with a
        capacity $Q_r$ ($r = 1, \dots, R$);
  \item for each job~$i$ and resource~$r$, the
        \emph{resource demand} $q_{ir}$: how many units of
        resource~$r$ job~$i$ uses while executing.
\end{itemize}
Resources are \emph{renewable}: a unit is occupied while a job
runs, then becomes available again.%
\aside{Workers, machines: busy during the
job, free afterwards.  Non-renewable
resources are consumed permanently.}

\medskip\noindent
A schedule must satisfy:
\begin{enumerate}[nosep]
  \item all precedence constraints (as before);
  \item at every time instant~$t$, the total demand for each
        resource~$r$ from all jobs running at time~$t$ does not
        exceed $Q_r$.
\end{enumerate}
The objective is to minimise the makespan $\Cmax$.
\end{definition}

The resource constraints are what make RCPSP fundamentally harder
than basic project scheduling.  Without resources, CPM schedules
every job as early as possible---there is no conflict.  With
resources, two jobs may both be ready to start, but starting
them simultaneously would exceed a resource capacity.  We must
choose which one goes first, and that choice affects all
subsequent decisions.

\begin{example}[RCPSP with two resources]
\label{ex:rcpsp-seven}
Take the seven-job instance from \cref{ex:seven-jobs} and add
two resources with capacities $Q_1 = 4$, $Q_2 = 3$:

\begin{figure}[H]
\centering
\begin{tabular}{lccccccc}
\toprule
\textbf{Job}      & 1 & 2 & 3 & 4 & 5 & 6 & 7 \\
\midrule
Duration $d_i$     & 3 & 3 & 1 & 4 & 2 & 1 & 3 \\
$q_{i,1}$ (R1)     & 3 & 2 & 3 & 1 & 0 & 2 & 1 \\
$q_{i,2}$ (R2)     & 0 & 2 & 1 & 2 & 2 & 1 & 2 \\
\bottomrule
\end{tabular}
\caption{Resource demands for the seven-job RCPSP instance.}
\label{fig:rcpsp-data}
\end{figure}

\noindent
With CPM (ignoring resources), jobs~2 and~3 start simultaneously
at time~3.  But job~2 needs 2~units of R1 and job~3 needs 3:
together that is $2 + 3 = 5 > 4 = Q_1$.  They cannot run in
parallel---we must choose.  This is a \textbf{choice point} that
does not exist in the unconstrained version.
\end{example}

% ------------------------------------------------------------
\subsection{Why RCPSP is hard}
\label{ssec:rcpsp-hard}
% ------------------------------------------------------------

Each choice point roughly doubles the number of candidate
schedules.  With $n$~jobs, the search space grows to approximately
$2^n$ alternatives.%
\aside{$2^n$ is a rough upper bound; the
true number depends on the graph
and resource constraints.}
For $n = 100$, that is $2^{100} \approx 10^{30}$
schedules---impossible to enumerate.

RCPSP is $\NP$-hard.  The PSPlib benchmark library contains
instances with 30, 60, 90, and 120~jobs;
some 120-job instances remain \emph{unsolved to optimality}
to this day.

% ------------------------------------------------------------
\subsection{Lower bound for RCPSP}
\label{ssec:rcpsp-lower-bound}
% ------------------------------------------------------------

A natural lower bound: \textbf{ignore the resource constraints}
and solve the resulting basic project scheduling problem with CPM\@.
Since removing constraints can only improve (or maintain) the
optimal value:
\[
  \Cmax^{\text{CPM}} \;\le\; \Cmax^{\text{RCPSP}*}.
\]
If a heuristic solution happens to match this bound, we know it
is optimal.  Otherwise, the gap
$\Cmax^{\text{heuristic}} - \Cmax^{\text{CPM}}$
tells us how far we \emph{might} be from optimality (though the
true optimum could lie anywhere in between).

% ------------------------------------------------------------
\subsection{Greedy heuristics for RCPSP}
\label{ssec:rcpsp-greedy}
% ------------------------------------------------------------

A greedy heuristic for RCPSP extends the CPM idea:
schedule one job at a time, choosing the \emph{earliest feasible}
start that respects both precedences and resources.
The question is: \textbf{which job to schedule next?}

There are two main families of greedy rules:

\paragraph{Time-based (schedule generation scheme).}
Advance time instant by instant.  At each moment, identify all
jobs that \emph{could} start (predecessors done, resources
available) and pick one by a priority rule:
longest processing time, shortest processing time, most resource
demand, etc.

\paragraph{Job-based (priority rule).}
Maintain three sets:
\begin{itemize}[nosep]
  \item \emph{Scheduled}: jobs already placed.
  \item \emph{Ready}: jobs whose predecessors are all scheduled.
  \item \emph{Not ready}: at least one predecessor is unscheduled.
\end{itemize}
At each step, select a job from the \emph{ready} set by a
priority rule, then find the earliest time at which it can start
without violating resource capacities.

\begin{remark}
Because the greedy traverses a single path through the exponential
search tree, different runs with different priority rules---or
even random tie-breaking---can produce very different makespans.
In the seven-job example, random restarts yielded solutions
ranging from $\Cmax = 14$ to $\Cmax = 16$, against a CPM lower
bound of~12.
\end{remark}

% ------------------------------------------------------------
\subsection{RCPSP extensions}
\label{ssec:rcpsp-extensions}
% ------------------------------------------------------------

Several extensions make the problem richer (and harder):

\begin{itemize}
  \item \textbf{Multi-mode RCPSP.}
    Each job has several \emph{execution modes}, each with its own
    duration and resource demands.
    For instance, digging a trench might take 3~days with
    3~workers, or 2~days with 1~worker and an excavator.
    The decision now includes both \emph{when} and \emph{how}
    to execute each job.

  \item \textbf{Non-renewable resources.}
    Resources that are \emph{consumed} (concrete, fuel, budget).
    In the single-mode case, total consumption is fixed regardless
    of schedule---feasibility reduces to checking whether the
    total demand exceeds supply.  Combined with multiple modes,
    however, the choice of mode affects consumption, making the
    problem significantly harder.

  \item \textbf{Preemption.}
    A job may be \emph{interrupted} and resumed later, freeing
    its resources for another job in the meantime.
    Whether preemption is realistic depends on the application:
    pausing a software build is easy; pausing a concrete pour
    is not.

  \item \textbf{Stochastic durations.}
    Durations drawn from probability distributions.
    Objectives become expected or worst-case makespan---this
    variant belongs more to the Project Management course.
\end{itemize}


% ============================================================
\section{Shop Scheduling Problems}
\label{sec:shop-scheduling}
% ============================================================

We now introduce a family of scheduling problems where
\emph{machines} (processors) are the central resource.
These are called \textbf{shop scheduling} problems.%
\aside{``Shop'' = workshop/factory floor,
not a retail store.}

The common setup: we have $n$ \emph{jobs} and $m$
\emph{machines}.  Each job consists of a sequence of
\emph{tasks} (also called \emph{operations}), where each task
must be processed on a specific machine for a given duration.
A machine can process at most one task at a time, and
(unless stated otherwise) a task cannot be interrupted once started.

% ------------------------------------------------------------
\subsection{Job Shop Scheduling (JSS)}
\label{ssec:job-shop}
% ------------------------------------------------------------

\begin{definition}[Job Shop Scheduling Problem]
\label{def:job-shop}
We are given:
\begin{itemize}[nosep]
  \item $n$ jobs, each consisting of an \emph{ordered} sequence
        of tasks;
  \item $m$ machines;
  \item each task $T_{ij}$ (the $j$-th task of job~$i$) is
        assigned to a specific machine $\mu_{ij}$ and has
        duration $d_{ij}$;
  \item (optionally) a \emph{release date} $r_i$ and a
        \emph{due date} $\bar{d}_i$ for each job.
\end{itemize}
\textbf{Constraints:}
\begin{enumerate}[nosep]
  \item \emph{Precedence within each job:} the tasks of a job
        must be processed in their given order.
  \item \emph{Machine capacity:} each machine processes at most
        one task at a time.
  \item Release dates are hard constraints.
\end{enumerate}
\textbf{Objective:} minimise the makespan $\Cmax$, or minimise
total \emph{tardiness} $\sum_i T_i$, where
$T_i = \max(0,\; C_i - \bar{d}_i)$ is the delay of job~$i$
past its due date.
\end{definition}

\begin{example}[A small job shop]
\label{ex:job-shop-small}
Three jobs on three machines:

\begin{figure}[H]
\centering
\begin{tikzpicture}[
    task/.style={draw=blue!60!black, fill=blue!8,
                 minimum height=7mm, font=\small,
                 rounded corners=1.5pt, anchor=west},
    arr/.style={-Stealth, thick},
    mlab/.style={font=\small\bfseries, anchor=east},
    x=8mm, y=10mm
  ]
  % Job 1: M3 -> M2 -> M1
  \node[mlab] at (-0.3, 3) {Job 1:};
  \node[task, fill=red!10, minimum width=32mm] (j1t1) at (0, 3)
    {\footnotesize $M_3$, $d=4$};
  \node[task, fill=red!10, minimum width=16mm] (j1t2) at (4.5, 3)
    {\footnotesize $M_2$, $d=2$};
  \node[task, fill=red!10, minimum width=8mm] (j1t3) at (6.8, 3)
    {\footnotesize $M_1$, $d=1$};
  \draw[arr] (j1t1) -- (j1t2);
  \draw[arr] (j1t2) -- (j1t3);

  % Job 2: M1 -> M3
  \node[mlab] at (-0.3, 2) {Job 2:};
  \node[task, fill=green!10, minimum width=24mm] (j2t1) at (0, 2)
    {\footnotesize $M_1$, $d=3$};
  \node[task, fill=green!10, minimum width=24mm] (j2t2) at (3.5, 2)
    {\footnotesize $M_3$, $d=3$};
  \draw[arr] (j2t1) -- (j2t2);

  % Job 3: M2 -> M1 -> M3
  \node[mlab] at (-0.3, 1) {Job 3:};
  \node[task, fill=yellow!15, minimum width=16mm] (j3t1) at (0, 1)
    {\footnotesize $M_2$, $d=2$};
  \node[task, fill=yellow!15, minimum width=32mm] (j3t2) at (2.5, 1)
    {\footnotesize $M_1$, $d=4$};
  \node[task, fill=yellow!15, minimum width=8mm] (j3t3) at (6.8, 1)
    {\footnotesize $M_3$, $d=1$};
  \draw[arr] (j3t1) -- (j3t2);
  \draw[arr] (j3t2) -- (j3t3);
\end{tikzpicture}
\caption{Three jobs in a job shop.  Each row shows one job's
         task sequence with its assigned machine and duration.
         Arrows indicate precedence within the job.}
\label{fig:job-shop-example}
\end{figure}

\noindent
Job~1 visits machines $3 \to 2 \to 1$; job~2 visits $1 \to 3$;
job~3 visits $2 \to 1 \to 3$.
The scheduler must decide, for each machine, in what order to
process the tasks assigned to it---without violating the
within-job precedences.
\end{example}

% ------------------------------------------------------------
\subsection{Flow Shop Scheduling (FSS)}
\label{ssec:flow-shop}
% ------------------------------------------------------------

\begin{definition}[Flow Shop Scheduling Problem]
\label{def:flow-shop}
A special case of the job shop where \textbf{all jobs visit the
machines in the same fixed order} $M_1 \to M_2 \to \cdots \to M_m$.
Every job has exactly $m$~tasks (some may have duration zero).
\end{definition}

The crucial simplification: since the machine order is
identical for every job, the \emph{only} decision is the
\textbf{permutation} (sequence) in which jobs enter the shop.%
\aside{If all machines use the same permutation, we have a
\emph{permutation flow shop}---the most common variant.}
Once the permutation is fixed, start times follow
deterministically: each task starts as soon as both its machine
is free and the previous task of the same job is finished.

\begin{figure}[H]
\centering
\begin{tikzpicture}[
    x=5mm, y=8mm,
    block/.style={draw, minimum height=6mm,
                  rounded corners=1pt, anchor=west,
                  font=\footnotesize},
    mlab/.style={font=\small\bfseries, anchor=east}
  ]
  % Machine labels
  \node[mlab] at (-0.5, 3) {$M_1$};
  \node[mlab] at (-0.5, 2) {$M_2$};
  \node[mlab] at (-0.5, 1) {$M_3$};

  % Job 3 (first in permutation)
  \node[block, fill=yellow!20, minimum width=15mm]
    (j3m1) at (0, 3) {$J_3$};
  \node[block, fill=yellow!20, minimum width=10mm]
    (j3m2) at (3, 2) {$J_3$};
  \node[block, fill=yellow!20, minimum width=10mm]
    (j3m3) at (4, 1) {$J_3$};

  % Job 1 (second)
  \node[block, fill=red!15, minimum width=10mm]
    (j1m1) at (3, 3) {$J_1$};
  \node[block, fill=red!15, minimum width=15mm]
    (j1m2) at (5, 2) {$J_1$};
  \node[block, fill=red!15, minimum width=10mm]
    (j1m3) at (8, 1) {$J_1$};

  % Job 2 (third)
  \node[block, fill=green!15, minimum width=10mm]
    (j2m1) at (5, 3) {$J_2$};
  \node[block, fill=green!15, minimum width=5mm]
    (j2m2) at (8, 2) {$J_2$};
  \node[block, fill=green!15, minimum width=15mm]
    (j2m3) at (10, 1) {$J_2$};

  % idle gap annotation
  \draw[decorate, decoration={brace, amplitude=3pt, mirror},
        gray]
    (5, 0.6) -- (8, 0.6)
    node[midway, below=4pt, font=\scriptsize\itshape, text=gray]
    {$M_3$ idle};

  % time axis
  \draw[-Stealth, gray] (-0.5, 0.2) -- (13, 0.2)
    node[right, font=\footnotesize] {time};
\end{tikzpicture}
\caption{A flow shop Gantt chart for the permutation
         $J_3, J_1, J_2$.  Machine~$M_3$ sits idle between
         $J_3$ and~$J_1$ because it must wait for $J_1$ to
         finish on~$M_2$.  A different permutation might reduce
         this idle time.}
\label{fig:flow-shop-gantt}
\end{figure}

\noindent
The goal is to find the permutation that minimises idle gaps and
thus the makespan.  Despite this structural simplification, the
flow shop with $m \ge 3$ machines is $\NP$-hard.

% ------------------------------------------------------------
\subsection{Open Shop Scheduling (OSS)}
\label{ssec:open-shop}
% ------------------------------------------------------------

\begin{definition}[Open Shop Scheduling Problem]
\label{def:open-shop}
Like the job shop, but with \textbf{no fixed ordering} of tasks
within a job.  Each job still has one task per machine, but the
scheduler is free to choose in which order a job visits the
machines.
\end{definition}

In the open shop, the within-job precedences disappear entirely.
What ties a job's tasks together are its release date and due
date: all tasks of a job must start after its release and finish
before its deadline.  The scheduler must decide both the
\emph{order of tasks within each job} and the \emph{order of
jobs on each machine}.

\medskip
\begin{center}
\begin{tabular}{lccc}
\toprule
 & \textbf{Machine order} & \textbf{Decision variable}
 & \textbf{Complexity} \\
\midrule
Job shop  & job-specific, fixed
          & start time of each task
          & $\NP$-hard \\
Flow shop & same for all jobs
          & job permutation
          & $\NP$-hard ($m \ge 3$) \\
Open shop & free
          & task order + start times
          & $\NP$-hard ($m \ge 3$) \\
\bottomrule
\end{tabular}
\end{center}

\begin{intermezzo}{Make as a project scheduler}
When you type \texttt{make} to compile a C++ project, the build
tool reads the dependency graph from the Makefile and solves a
project scheduling problem: each compilation command is a job,
its duration is the compilation time, and ``\texttt{A.o} depends
on \texttt{A.cc}'' is a precedence constraint.

With \texttt{make -j} (parallel mode), Make schedules independent
compilations on separate CPU cores---exactly the CPM forward
strategy.  In a classroom demo, a synthetic Makefile encoding the
seven-job instance from \cref{ex:seven-jobs} (using \texttt{sleep}
to simulate durations) completed in 12~seconds with
\texttt{make -j}, matching the CPM makespan, versus 16~seconds
in sequential mode.
\end{intermezzo}
